\centerline{\bf \Huge Геометрия I}
\centerline{\bf \Huge Признаки равенства}

\begin{flushright}
        \scriptsize{Люди живут только потому, что могут забывать прошлое.\\ Но есть вещи, которые нужно помнить всегда.\\Гендо Икари}
\end{flushright}

\problem{1}
Равные отрезки $A B$ и $C D$ пересекаются в точке $O$, причём $A O=O D$. Докажите равенство треугольников $A B C$ и $D C B$.

\problem{2}
Точки $A, B, C, D$ лежат на одной прямой. Докажите, что если треугольники $A B E_1$ и $A B E_2$ равны, то треугольники $C D E_1$ и $C D E_2$ тоже равны.

\problem{3}
Высоты треугольника $A B C$, проведённые из вершин $B$ и $C$ пересекаются в точке $M$. Известно, что $B M=C M$. Докажите, что треугольник $A B C-$ равнобедренный.

\problem{4}
В треугольнике $A B C$ проведена биссектриса $B L$. Известно, что $B L=A B$. На продолжении $B L$ за точку $L$ выбрана точка $K$, причём $\angle B A K+\angle B A L=180^{\circ}$. Докажите, что $B K=B C$.

\problem{5}
Через вершину $A$ квадрата $A B C D$ проведены прямые $l_1$ и $l_2$, пересекающие его стороны. Из точек $B$ и $D$ опущены перпендикуляры $B B_1, B B_2, D D_1$ и $D D_2$ на эти прямые. Докажите, что отрезки $B_1 B_2$ и $D_1 D_2$ равны и перпендикулярны.

\problem{6}
На сторонах $A B$ и $B C$ равностороннего треугольника $A B C$ отмечены точки $D$ и $K$ соответственно, а на стороне $A C$ отмечены точки $E$ и $M$ так, что $D A+A E=K C+C M=A B$. Отрезки $D M$ и $K E$ пересекаются. Найдите угол между ними.